\paragraph{窗口分辨率压缩的不可压缩余量：有限账本、半尺度零块与最小退化扇区}\label{par:gut-resolution-ledger-collision-obstruction}
在结论 \ref{par:gut-foldbin-fold-complement-identity}--\ref{par:gut-d18-endogenous-split} 的同层补余链之外，还可抽取四条更偏“不可压缩性”的刚性后果：
有限层异常账本不能在严格乘法到加法同态语义下承载算术；
六倍窗零块至少覆盖一个半尺度完整稳定型空间；
在 bin-fold 退化谱口径下，只要最小退化 Fibonacci 律沿偶窗延拓，最小退化扇区单独就足以驱动二阶碰撞尺度发散；
并且在已审计偶窗上，总预算 \(16\) 的最小退化包只在 \(m=6\) 出现。

\begin{theorem}[有限分辨率异常账本不可能同态承载乘法算术]\label{thm:gut-finite-anomaly-ledger-no-multiplicative-embedding}
令
\[
\Lambda_{\mathrm{anom}}
:=\ZZ^a\oplus\bigoplus_{j=1}^r\bigl(\mathcal G^{\mathrm{bin}}_{m_j}\bigr)^{\mathrm{ab}},
\qquad a,r<\infty,\quad m_j\ge 1.
\]
则不存在单射幺半群同态
\[
\Phi:(\NN_{>0},\times)\hookrightarrow (\Lambda_{\mathrm{anom}},+).
\]
特别地，任何有限层的 fold-gauge 异常寄存器直和，都不可能在保持“乘法 \(\to\) 加法”严格同态语义的前提下充当素因子账本。
在 \(m=6\) 时，由定理 \ref{thm:window6-foldbin-gauge-abelianization-even-parity}，
\[
\bigl(\mathcal G^{\mathrm{bin}}_6\bigr)^{\mathrm{ab}}\cong (\ZZ_2)^{21},
\]
故 window-\(6\) 的完整可见符号电荷格亦属于该禁例。
\leanverified{paper\_gut\_finite\_anomaly\_no\_multiplicative\_embedding\_seeds}
\leanverified{paper\_gut\_finite\_anomaly\_no\_multiplicative\_embedding\_package}
\leanverified{paper\_gut\_finite\_anomaly\_ledger\_no\_multiplicative\_embedding}
\end{theorem}
\begin{proof}
由定理 \ref{thm:window6-foldbin-gauge-abelianization-even-parity}，每个
\(
\bigl(\mathcal G^{\mathrm{bin}}_{m_j}\bigr)^{\mathrm{ab}}
\)
都是有限 \(2\)-群。
因此
\[
T:=\bigoplus_{j=1}^r\bigl(\mathcal G^{\mathrm{bin}}_{m_j}\bigr)^{\mathrm{ab}}
\]
为有限阿贝尔群。
记其指数为
\(
e:=\exp(T)\ge 1
\)；
则对任意 \(y\in T\) 有 \(ey=0\)。
\par
反设存在单射幺半群同态 \(\Phi\)。
任取两两不同的素数 \(p_1,\dots,p_{a+1}\)。
由于 \(\Phi(n^e)=e\,\Phi(n)\) 且 \(e\) 杀死挠部分，故
\[
\Phi(p_i^e)\in \ZZ^a\oplus\{0\}\cong \ZZ^a
\qquad(1\le i\le a+1).
\]
于是 \(a+1\) 个向量 \(\Phi(p_i^e)\) 在 \(a\) 维有理向量空间 \(\QQ^a\) 中线性相关。
清分母可得存在整数 \(c_1,\dots,c_{a+1}\)，不全为零，使
\[
\sum_{i=1}^{a+1} c_i\,\Phi(p_i^e)=0.
\]
将 \(c_i\) 分解为
\(
c_i=c_i^+-c_i^-
\)
（其中 \(c_i^\pm\in\NN\)，且 \(c_i^+c_i^-=0\)），则
\[
\sum_{i=1}^{a+1} c_i^+\,\Phi(p_i^e)
=
\sum_{i=1}^{a+1} c_i^-\,\Phi(p_i^e).
\]
由 \(\Phi\) 的幺半群同态性，
\[
\Phi\!\left(\prod_{i=1}^{a+1} p_i^{e c_i^+}\right)
=
\Phi\!\left(\prod_{i=1}^{a+1} p_i^{e c_i^-}\right).
\]
由于 \(\Phi\) 单射，必有
\[
\prod_{i=1}^{a+1} p_i^{e c_i^+}
=
\prod_{i=1}^{a+1} p_i^{e c_i^-}.
\]
算术基本定理遂强制 \(c_i^+=c_i^-\) 对一切 \(i\) 成立，从而 \(c_i=0\) 全部成立，与所取关系非零矛盾。
故此类 \(\Phi\) 不存在。证毕。
\end{proof}

\leanverified{paper\_gut\_half\_scale\_zero\_block\_covers\_stable\_space}
\begin{theorem}[六倍窗零块至少覆盖一个半尺度完整稳定型空间]\label{thm:gut-half-scale-zero-block-covers-stable-space}
若
\[
m+2\equiv 0\pmod 6,
\]
则谱零点集合 \(\mathcal Z_m\) 满足
\[
\boxed{\ |\mathcal Z_m|\ \ge\ F_{(m+2)/2}=|X_{(m-2)/2}|\ }.
\]
换言之，六倍窗上的半阶零块至少覆盖一个半尺度完整稳定型空间的基数。
\end{theorem}
\begin{proof}
由结论 \ref{con:xi-fold-window6-kernel-half-index-lower-bound}，
\[
|\mathcal Z_m|\ge F_{(m+2)/2}.
\]
另一方面，稳定型集合满足标准 Fibonacci 计数律
\[
|X_n|=F_{n+2}\qquad(n\ge 0).
\]
取
\(
n=(m-2)/2
\)
即得
\[
|X_{(m-2)/2}|=F_{(m+2)/2}.
\]
合并两式即得结论。证毕。
\end{proof}

\begin{remark}[大窗数值核验]
由脚本 \texttt{\detokenize{scripts/exp_fold_zero_coset_union_count.py}} 的精确并集计数，
在 \(m=58\) 时有
\[
|\mathcal Z_{58}|=839415>F_{30}=832040=|X_{28}|.
\]
该值与定理 \ref{thm:gut-half-scale-zero-block-covers-stable-space} 的半尺度下界完全相容，并显示该下界在大窗上仍保留可见余量。
\end{remark}

\begin{theorem}[全偶窗 Fibonacci 延拓假设下最小退化扇区单独驱动 bin-fold 二阶碰撞发散]\label{thm:gut-bin-min-sector-forces-collision-divergence}
定义 bin-fold 二阶碰撞量
\[
\mathsf{Col}^{\mathrm{bin}}_m
:=2^{-2m}\sum_{x\in X_m}\bigl(d^{\mathrm{bin}}_m(x)\bigr)^2.
\]
设存在偶整数 \(m_0\) 使得对每个偶数 \(m\ge m_0\) 都有
\[
d_{\min}(m)=F_{m/2},
\qquad
|\mathcal S_{m,d_{\min}(m)}|=F_m.
\]
则对所有此类 \(m\) 成立显式下界
\[
F_{m+2}\,\mathsf{Col}^{\mathrm{bin}}_m
\ \ge\
\frac{F_{m+2}F_mF_{m/2}^2}{4^m}.
\]
并且当 \(m\to\infty\) 沿偶整数时，
\[
\frac{F_{m+2}F_mF_{m/2}^2}{4^m}
=
\frac{\varphi^2}{25}\Bigl(\frac{\varphi^3}{4}\Bigr)^m(1+o(1)).
\]
特别地，
\[
F_{m+2}\,\mathsf{Col}^{\mathrm{bin}}_m\longrightarrow +\infty.
\]
因此在该延拓假设下，无需求助全部退化谱，最小退化扇区一项已经足以驱动 bin-fold 二阶碰撞尺度发散。
\leanverified{paper\_gut\_bin\_min\_sector\_forces\_collision\_divergence}
\end{theorem}
\begin{proof}
由 \(\mathsf{Col}^{\mathrm{bin}}_m\) 的定义，对任意子集 \(A\subseteq X_m\) 都有
\[
\mathsf{Col}^{\mathrm{bin}}_m
\ge
2^{-2m}\sum_{x\in A}\bigl(d^{\mathrm{bin}}_m(x)\bigr)^2.
\]
取 \(A=\mathcal S_{m,d_{\min}(m)}\)，并用该扇区上
\(
d^{\mathrm{bin}}_m(x)=d_{\min}(m)
\)
的定义，得到
\[
\mathsf{Col}^{\mathrm{bin}}_m
\ge
2^{-2m}\,|\mathcal S_{m,d_{\min}(m)}|\,d_{\min}(m)^2.
\]
在假设下代入
\(
|\mathcal S_{m,d_{\min}(m)}|=F_m
\)
与
\(
d_{\min}(m)=F_{m/2}
\)，
并再乘以 \(F_{m+2}\)，即得
\[
F_{m+2}\,\mathsf{Col}^{\mathrm{bin}}_m
\ge
\frac{F_{m+2}F_mF_{m/2}^2}{4^m}.
\]
再由 Binet 渐近
\[
F_k=\frac{\varphi^k}{\sqrt5}\,(1+o(1))
\qquad(k\to\infty)
\]
得到
\[
\frac{F_{m+2}F_mF_{m/2}^2}{4^m}
=
\frac{\varphi^{m+2}}{\sqrt5}\cdot
\frac{\varphi^m}{\sqrt5}\cdot
\frac{\varphi^m}{5}\cdot
4^{-m}\,(1+o(1))
=
\frac{\varphi^2}{25}\Bigl(\frac{\varphi^3}{4}\Bigr)^m(1+o(1)).
\]
由于 \(\varphi^3/4>1\)，右端发散到 \(+\infty\)，故结论成立。证毕。
\end{proof}

\begin{theorem}[已审计偶窗中总预算 \texorpdfstring{$16$}{16} 的最小退化包唯一性]\label{thm:gut-audited-even-windows-unique-budget16-min-sector}
在已审计偶窗
\[
m\in\{6,8,10,12\}
\]
上，定义最小退化包的总微态预算
\[
B_m:=d_{\min}(m)\,|\mathcal S_{m,d_{\min}(m)}|.
\]
则由表 \ref{tab:fold_bin_degeneracy_spectrum_m6_m12} 有
\[
d_{\min}(m)=F_{m/2},
\qquad
|\mathcal S_{m,d_{\min}(m)}|=F_m
\qquad
\bigl(m\in\{6,8,10,12\}\bigr),
\]
从而
\[
B_m=F_{m/2}F_m
\qquad
\bigl(m\in\{6,8,10,12\}\bigr),
\]
并显式给出
\[
(B_6,B_8,B_{10},B_{12})=(16,63,275,1152).
\]
因此在已审计偶窗中，
\[
\boxed{\ B_m=16\iff m=6\ }.
\]
此外，window-\(6\) 同时满足
\[
|\mathcal S_{6,d_{\min}(6)}|=8,
\qquad
B_6=16,
\qquad
|X_6\setminus \mathcal S_{6,d_{\min}(6)}|=D_{10}=13.
\]
\leanverified{paper\_gut\_budget16\_unique\_m6}
\leanverified{paper\_boundary\_tower\_fib\_count}
\end{theorem}
\begin{proof}
表 \ref{tab:fold_bin_degeneracy_spectrum_m6_m12} 直接给出
\[
d_{\min}(6)=2,\quad d_{\min}(8)=3,\quad d_{\min}(10)=5,\quad d_{\min}(12)=8,
\]
\[
|\mathcal S_{6,2}|=8,\quad |\mathcal S_{8,3}|=21,\quad |\mathcal S_{10,5}|=55,\quad |\mathcal S_{12,8}|=144.
\]
这正好分别等于
\[
(F_3,F_4,F_5,F_6),\qquad (F_6,F_8,F_{10},F_{12}),
\]
故
\[
B_6=F_3F_6=2\cdot 8=16,\qquad
B_8=F_4F_8=3\cdot 21=63,
\]
\[
B_{10}=F_5F_{10}=5\cdot 55=275,\qquad
B_{12}=F_6F_{12}=8\cdot 144=1152.
\]
于是 \(B_m\) 在该四个审计点上严格递增，从而预算 \(16\) 仅在 \(m=6\) 出现。
最后，由接口化结论 \ref{par:gut-foldbin-fold-complement-identity} 在 \(m=6\) 的特例，
\[
|X_6\setminus \mathcal S_{6,d_{\min}(6)}|=D_{10}=13.
\]
证毕。
\leanverified{fib\_product\_strict\_mono}
\leanverified{budget\_strict\_mono}
\leanverified{paper\_gut\_budget16\_strict\_mono\_extended}
\end{proof}
